% Generated draft: requires mathematical and source review.
\documentclass[11pt]{article}
\usepackage[margin=25mm]{geometry}
\usepackage{fontspec}
\setmainfont{Latin Modern Roman}
\usepackage{amsmath,amssymb,mathtools,unicode-math}
\setmathfont{Latin Modern Math}
\usepackage{enumitem,needspace,xurl,hyperref}
\setlength{\parindent}{0pt}
\setlength{\parskip}{5pt}
\setlength{\emergencystretch}{3em}
\newcommand{\N}{\mathbb{N}}
\newcommand{\Z}{\mathbb{Z}}
\newcommand{\Q}{\mathbb{Q}}
\newcommand{\R}{\mathbb{R}}
\newcommand{\C}{\mathbb{C}}
\newcommand{\F}{\mathbb{F}}
\newcommand{\A}{\mathbb{A}}
\newcommand{\PP}{\mathbb{P}}
\newcommand{\E}{\mathbb{E}}
\newcommand{\eps}{\varepsilon}
\newcommand{\dd}{\,\mathrm{d}}
\begin{document}
\section*{UnsolvedMath ID: 30001419}
\textbf{Source problem code:} catalogue entry
\subsection{Definitions and mathematical statement}
A thin-coating scattering problem is understood here only descriptively as a scattering problem involving a coating of small thickness. A generalized impedance boundary condition is an effective boundary condition intended to represent the coating's influence on the scattering field. The supplied record does not specify the governing equations, the coating geometry or material parameters, or the boundary operators involved. Its phrase ``boundary contains an angle'' indicates an angular boundary feature, but does not specify its geometric meaning or admissible angles.

The supplied research question is: for thin-coating scattering problems whose boundary contains an angle, can high-order generalized impedance boundary conditions be derived, and can the resulting asymptotic models be analyzed rigorously in the time domain and in more general settings?

This question does not determine a precise quantified mathematical assertion. In particular, the spatial dimension, scattering equations, interface and exterior conditions, thickness regime, material assumptions, requested asymptotic orders, treatment of the angular feature, and meaning of ``more general settings'' remain unspecified. For a time-domain formulation, initial data, temporal assumptions, and compatibility conditions also remain unspecified. No particular well-posedness statement, error estimate, norm, or interval of validity is supplied. These choices must be established from reliable source context before a precise theorem or conjecture can be formulated.
\subsection{Short English statement}
The question asks whether high-order effective impedance conditions can describe scattering by thin coatings near angular boundary features, with rigorous justification in the time domain and broader settings. The supplied record leaves the equations, hypotheses, approximation orders, and required estimates unspecified, so it supports a research question rather than a precise conjecture.
\subsection{Sources}
\begin{itemize}
\item Local UnsolvedMath registry, supplied record with permanent ID \(30001419\), source code \texttt{catalogue entry}, titled ``High-Order Impedance Conditions Near Boundary Corners,'' in the Oberwolfach Reports Open Problems collection. The proposed year is recorded as 2010; no individual proposer is supplied.
\item \emph{Computational Electromagnetism and Acoustics} (2010), Oberwolfach Reports, report 10/2010, pp.~415--519, DOI \texttt{10.4171/owr/2010/10}. Supplied source and literature reference: \url{https://doi.org/10.4171/owr/2010/10}. The supplied background describes Oberwolfach Reports as records of mathematical workshops and their open-problem discussions.
\item The registry labels the statement ``exact'' and reports that its mathematical content was checked against the official report, with notation, context, and question form normalized without changing the source-level claim. This is supplied verification metadata, not an independent verification performed for this draft.
\item Literature assessment dated exactly \texttt{2026-08-21}: the supplied curation labels the record open, explaining that the extraction is only an open-questions heading without a specific mathematical claim, so resolution cannot be established and the open label follows its requested classification rule. Its listed checked source is Oberwolfach Report 10/2010 at \url{https://doi.org/10.4171/owr/2010/10}. This assessment is supplied dated curation, not independently verified; its time-sensitive status and cited claims require independent checking.
\end{itemize}
\end{document}
